\documentclass[11pt,a4paper]{article}
% ---------------------------------------------------------------------------
%  Summa -- the style of the showcase and of the manual: the common style of the packages, n4sd4q.sty (a copy of
%  the one of the hub, with the palette n4sd4q-palette.tex: found through TEXINPUTS, set by build_pdf.sh and
%  build_manual.sh), and what is particular to Summa. The examples are computed by Development/run_examples.wl:
%  examples/NAME.wl is the code, NAME.tex the output as a formula, NAME.txt the output as text, NAME.pdf the output
%  as a picture (a diagram, a contraction tree, a plot, the summary box of an EinsumFunction).
% ---------------------------------------------------------------------------
\usepackage[utf8]{inputenc}
\usepackage{n4sd4q}
\familysetup{package=Summa, logo=../Logo/Summa_logo_white.pdf}

\familykeywords{Einsum,EinsumPath,EinsumFunction,EinsumDiagram,EinsumTree,EinsumClearCache,SummaUpdateCheck}
\lstset{literate={α}{{$\alpha$}}1 {β}{{$\beta$}}1 {γ}{{$\gamma$}}1 {δ}{{$\delta$}}1 {ε}{{$\varepsilon$}}1 {ϵ}{{$\epsilon$}}1 {η}{{$\eta$}}1
  {μ}{{$\mu$}}1 {ν}{{$\nu$}}1 {σ}{{$\sigma$}}1 {τ}{{$\tau$}}1 {ψ}{{$\psi$}}1 {φ}{{$\varphi$}}1 {ϕ}{{$\phi$}}1 {χ}{{$\chi$}}1
  {ω}{{$\omega$}}1 {→}{{$\to$}}1 {…}{{$\ldots$}}1 {·}{{$\cdot$}}1 {∑}{{$\sum$}}1 {×}{{$\times$}}1}
% the same letters in the text (\texttt{α[k]}, as the code writes them)
\DeclareUnicodeCharacter{03B1}{\ensuremath{\alpha}}\DeclareUnicodeCharacter{03B2}{\ensuremath{\beta}}
\DeclareUnicodeCharacter{03C3}{\ensuremath{\sigma}}\DeclareUnicodeCharacter{2192}{\ensuremath{\to}}

\newcommand{\exmissing}[1]{\textcolor{red}{[missing example: \detokenize{#1}]}}
\makeatletter
\newcommand{\excodeafter}{\tcb@parfillskip@check\par\ifvmode\nobreak\@nobreaktrue\nointerlineskip\addvspace{\glueexpr2pt-\parskip}\fi}
\makeatother
\newcommand{\excode}[1]{\ifhmode{\widowpenalty=10000\par}\nobreak\fi\IfFileExists{examples/#1.wl}{\tcbinputlisting{listing only,
  colback=codebg, colframe=codeframe, boxrule=0.5pt, arc=2pt, left=6pt, right=6pt, top=3pt, bottom=3pt, before skip=4pt,
  after=\excodeafter, listing file={examples/#1.wl}, listing options={language=WL}}}{\exmissing{#1}}}
\newcommand{\exformula}[1]{\begin{exoutput}\raggedright\IfFileExists{examples/#1.tex}{$\displaystyle\input{examples/#1.tex}$}{\exmissing{#1}}\end{exoutput}}
\newcommand{\extext}[1]{\IfFileExists{examples/#1.txt}{\tcbinputlisting{listing only, colback=white, colframe=codeframe,
  boxrule=0.5pt, arc=2pt, left=24pt, right=6pt, top=2pt, bottom=2pt, before skip=2pt, after skip=8pt,
  listing file={examples/#1.txt}, listing options={style=plainoutput}, overlay unbroken and first={\node[anchor=north west,
  font=\scriptsize\sffamily\color{comment}] at ([xshift=4pt, yshift=-5pt]frame.north west) {out};}}}{\exmissing{#1}}}
% a picture, centred in the box of the output, on white: at the size of its PDF (the front end exports it at 72/100 of
% its size on the screen), or as wide as the box if it is wider (\exprepare{width}{NAME} sets the width instead). The code
% and the picture stay on the same page
\newsavebox{\expicturebox}
\newcommand{\exprepare}[2]{\ifx\relax#1\relax\sbox{\expicturebox}{\includegraphics{examples/#2.pdf}}%
  \ifdim\wd\expicturebox>\dimexpr\linewidth-38pt\relax
  \sbox{\expicturebox}{\resizebox{\dimexpr\linewidth-38pt\relax}{!}{\includegraphics{examples/#2.pdf}}}\fi
  \else\sbox{\expicturebox}{\includegraphics[width=#1]{examples/#2.pdf}}\fi}
\newcommand{\expicture}{\begin{exoutput}\centering\usebox{\expicturebox}\end{exoutput}}
% room for a picture and its code: a new page when the rest of this one is too short. Unlike \needspace, no negative
% penalty is left when the room is there, so that the page is not broken here anyway; \keepfor[lines]{NAME} before the
% text that leads to a picture (headings included, about so many lines with the code) keeps the text with it
\makeatletter
\newcommand{\exkeepspace}[1]{\par\penalty\@M\begingroup\dimen@=\dimexpr#1\relax\dimen@ii=\pagegoal\advance\dimen@ii-\pagetotal
  \ifdim\dimen@>\dimen@ii\ifdim\dimen@ii>\z@\vfil\penalty-\@M\fi\fi\endgroup}
\makeatother
\newcommand{\expicturekeep}[2]{\IfFileExists{examples/#2.pdf}{\exprepare{#1}{#2}%
  \exkeepspace{\ht\expicturebox+\dp\expicturebox+8\baselineskip}}{}}
\newcommand{\keeplines}[1]{\exkeepspace{#1\baselineskip}}   % room for a table and the sentence before it
\newcommand{\keepfor}[2][8]{\IfFileExists{examples/#2.pdf}{\exprepare{}{#2}%
  \exkeepspace{\ht\expicturebox+\dp\expicturebox+#1\baselineskip+3\baselineskip}}{}}   % [lines of text and code before it]
\newcommand{\example}[1]{\expicturekeep{}{#1}\excode{#1}\IfFileExists{examples/#1.tex}{\exformula{#1}}{\IfFileExists{examples/#1.txt}{\extext{#1}}%
  {\IfFileExists{examples/#1.pdf}{\expicture}{}}}}
% an example whose picture is drawn at a given width
\newcommand{\examplewidth}[2]{\expicturekeep{#1}{#2}\excode{#2}\IfFileExists{examples/#2.pdf}{\expicture}{\exmissing{#2}}}

\familysetup{document=Showcase}

\begin{document}
\familytitle{Tensor contractions written as on paper: the package finds the order and the kernels}{Everything the package
does, step by step: every output is computed by the code printed above it}{version 1.0.0 \textperiodcentered{} November 2026}

\begin{tcolorbox}[colback=accentpale, colframe=accentpale, arc=2pt, left=8pt, right=8pt]
\small Summa is a Wolfram Language package for tensor contractions, written with the syntax of \texttt{numpy.einsum}. A
contraction is written as on paper, by the indices of every factor: \texttt{"ij,jk->ik"} is $\sum_ja_{ij}b_{jk}$. How to
compute it is left to the package: it chooses the order of the pairwise contractions, which in a network of tensors can
change the cost by orders of magnitude, and the kernel of each, BLAS or a vectorised one, with no transposes, reshapes or
intermediate arrays written by hand. An index can appear any number of times, in one tensor or in many. This document goes
through all the package does, from a matrix product to tensor networks and a Kubo formula: the syntax; indices repeated at
will; the order of the contractions and the pictures of a network; compiled einsums for loops; symbolic, exact, complex
and sparse arrays; computations of physics; how a contraction runs and how the package is checked. The code of each
section uses only what the previous ones introduced.
\end{tcolorbox}

{\small\tableofcontents}

% =====================================================================================================================
\section{Written as on paper}

\subsection{One string for the whole contraction}
The package is loaded with
\example{load}
Each operand gets a subscript, a string of one-character indices; the subscripts are separated by commas, and the arrow is
followed by the indices of the result. An index that appears in the operands and not in the result is summed over:
\texttt{Einsum["ij,jk->ik", m, n]} is $\sum_jm_{ij}n_{jk}$, the matrix product,
\example{first}
\example{firstcheck}
How it is computed, here by one call of \texttt{Dot}, is the business of the package, for this contraction and for every
one below. One operand is enough for a transpose, a trace, a diagonal and sums over rows or over all entries:
\example{onematrix}
Two vectors give the inner and the outer product:
\example{vectors}

\subsection{The implicit output}
Without the arrow, the result has the indices that appear exactly once, in alphabetical order, as in numpy:
\texttt{"ij,jk"} is the matrix product, \texttt{"ba"} the transpose and \texttt{"ii"} the trace:
\example{implicit}

\subsection{Diagonals}
An index repeated within one operand takes a diagonal: \texttt{"iij->ij"} is $t_{iij}$, and \texttt{"iii->i"} the
diagonal of a cube. The entries of \texttt{digits} are $100i+10j+k$, so the results show which ones were taken:
\example{diagonals}

\subsection{The ellipsis and broadcasting}
\texttt{...} stands for the leading indices that are not named, as in numpy; here a stack of five $2\times3$ matrices
times one $3\times4$ matrix:
\example{ellipsis}
Axes of size 1 are broadcast: the single row of the second operand is paired with every row of the first:
\example{broadcast}

\subsection{Scalars}
An empty subscript is a scalar operand. Any argument that is not a list or a \texttt{SparseArray}, \texttt{Sqrt[2]} or
\texttt{a + b} as well, is a scalar:
\example{scalars}

\subsection{Labels of any kind}
In the list syntax the subscripts are lists, and an index can be any expression: a string, an integer, a symbol with
arguments. The string syntax also takes Unicode letters, spaces and the arrow $\to$:
\example{listsyntax}
\example{unicode}
Indices with arguments, $\alpha_k$ and $\sigma_k$ written \texttt{α[k]} and \texttt{σ[k]}, let a program write the
subscripts of a network of any size (section~\ref{sec:mps}).

% =====================================================================================================================
\section{Indices repeated at will}\label{sec:repeated}
On paper an index appears in as many factors as needed: $\sum_ix_iy_iz_i$, or a batch index $b$ that labels independent
products and stays in the result, as in $\sum_jL_{bij}R_{bjk}$. Summa allows an index any number of times, within one tensor
and across tensors, and this is the reason it exists. The two tools of the Wolfram Language for the job stop short of
it. \texttt{ResourceFunction["EinsteinSummation"]} refuses an index that appears three times
(\texttt{\{\{i\}, \{i\}, \{i\}\} -> \{\}}) and a batch index kept in the result
(\texttt{\{\{k, a, b\}, \{k, b, c\}\} -> \{k, a, c\}}). \texttt{TensorContract} can sum over three or more slots, but of the
outer product of all the operands, which it builds first, and it cannot keep a shared index in the result (a batch index,
a diagonal).

\subsection{One index in three tensors}
\example{threevectors}
\texttt{TensorContract} gives the same sum from the outer product:
\example{tensorcontract}
In the diagrams of section~\ref{sec:pictures} such an index is a hyperedge, a dot joined to all the tensors that share it.

\subsection{Batch indices}
An index shared by the operands and kept in the result labels independent contractions, all done at once. The quadratic
form $\sum_{ij}x_{bi}q_{ij}x_{bj}$ of every row $b$, with $b$ in two operands and in the result:
\example{quadraticforms}
A thousand matrix products in one call:
\example{batchmatmul}

\subsection{The CP reconstruction}
A tensor in the canonical polyadic form is a sum of outer products, $T_{abc}=\sum_iA_{ai}B_{bi}C_{ci}$: the index $i$ is in
three tensors, and $a$, $b$, $c$ are in the result. With $r=8$ terms:
\example{cp}
\texttt{TensorContract} would first build the outer product of the three matrices, $40^3\cdot8^3$ numbers. Summa contracts
two of them, keeping $i$ as a batch index, and then the third, with half a million multiplications (the cost is given by
\fn{EinsumPath}, section~\ref{sec:order}):
\example{cpcost}
On the author's computer, with Mathematica 15, the outer product takes 262\,MB and 0.37\,s, the einsum 1.6\,ms.

\subsection{Shared and kept}
An index can also be shared and kept without any other: the dot products of the rows of two matrices,
$\sum_jm_{ij}n_{ij}$ for every $i$:
\example{rowdots}

% =====================================================================================================================
\section{The order of the contractions}\label{sec:order}

\subsection{What a path costs}
An einsum of $n$ tensors is computed as $n-1$ contractions of two tensors, and their order decides the cost.
\fn{EinsumPath} tells which order the package chooses and what it costs, from the shapes alone (the tensors or their
dimensions). For the product of three matrices, $1000\times10$, $10\times1000$ and $1000\times10$:
\example{chain}
\opt{Path} is in the linear format of numpy and opt\_einsum~\cite{SG18}, 1-based: each pair gives the positions of two
tensors in the current list; they are removed, and their result is appended at the end. Here $B$ and $C$ go first,
then $A$ with their product. The cost of a step is the number of terms of its sum, the product of the dimensions of all its
indices; \opt{Flops} is the total, \opt{NaiveFlops} the cost of the sum written out over all the indices at once,
\opt{Speedup} their ratio and \opt{LargestIntermediate} the size of the largest tensor built on the way. Taking $A$ times
$B$ first costs a hundred times more, and builds a $1000\times1000$ intermediate:
\example{chainbad}

\subsection{Strategies}
\texttt{"Path" -> "Optimal"} finds the cheapest order exactly, by dynamic programming over the subsets of the tensors. The
result of a subset keeps the indices still needed outside it or in the output, which is how an index shared by three
tensors becomes a batch index of the intermediates. It is the default up to 8 tensors and is accepted up to 14; beyond 8
the default is \opt{Greedy}, the search of opt\_einsum, which contracts at each step the pair that least increases the size.
\opt{Sequential} contracts from left to right, and a list of pairs is a path of your own. On a network of five tensors:
\example{strategies}

\subsection{The norm of a matrix product state}\label{sec:mps}
A matrix product state (MPS) of $L$ sites has a tensor $A^{(k)}_{\alpha_{k-1}\sigma_k\alpha_k}$ on each site, with a
physical index $\sigma_k$ of dimension $d$ and bonds $\alpha_k$ of dimension $\chi$. Its norm
$\langle\psi|\psi\rangle$ joins it to its conjugate over the physical indices. For four sites with $d=2$ and $\chi=64$ the
sum over all the indices at once has $2^4\cdot64^6\approx7\times10^{10}$ terms of eight factors; the order chosen by the
package costs $7\times10^4$ multiplications, a speed-up of $10^8$:
\example{mpsnorm}
The order matters also among the reasonable ones. The function \texttt{mps} writes the norm of an MPS of any length in the
list syntax, with the indices \texttt{α[k]} (bonds of the ket), \texttt{β[k]} (bonds of the bra) and \texttt{σ[k]}; the
tensors at the ends have one bond. \texttt{mpsShapes} gives their dimensions:
\example{mpsgeneral}
Contracting the ket first, in the order in which the tensors are written (\opt{Sequential}), builds the whole state, $d^L$
amplitudes: the cost grows exponentially with $L$. The order of the package keeps every intermediate small, as zipping
ket and bra together site by site does, and its cost grows linearly; at twenty sites it is more than two thousand times
cheaper:
\example{mpstable}

% =====================================================================================================================
\keepfor[10]{diagrammps}\section{Pictures of a network}\label{sec:pictures}

\subsection{Bonds, open legs, traces, sums and hyperedges}
\fn{EinsumDiagram} draws the network of an einsum: tensors are disks, a bond joins two tensors that share an index, and an
index of the result is an open leg, in blue. The MPS of four sites, with its physical legs open:
\example{diagrammps}
\keepfor[9]{diagramkinds}Every index is drawn as a COPY tensor (a delta) with one leg for each time it appears, plus one if it is in the result.
Two legs are a wire: a bond, an open leg, or a loop for a trace. One leg is an index summed in one tensor alone, drawn as
a dashed leg ending in a small circle (a contraction with the vector of ones). Three or more legs are a hyperedge, a purple
dot joined to all the tensors that share it. All of them in one network: the trace $i$ and the summed index $x$ on $A$,
the bonds $j$ and $k$, the hyperedge $h$ shared by $B$, $C$ and $D$, and the open leg $l$:
\example{diagramkinds}
Tensors with the same name have the same colour, and tensors passed as symbols are named after them. The layout is
computed by classical multidimensional scaling of the distances in the graph, refined by stress
majorization~\cite{GKN04}: chains come out straight, rings round and grids square.

\keepfor[10]{diagrampath}\subsection{Dimensions and the path}
Given the dimensions, the diagram writes them next to the indices and draws thicker bonds for larger dimensions (on a
logarithmic scale); \texttt{"ShowPath" -> True} marks each bond with the step at which it is contracted. The norm of the
MPS of section~\ref{sec:mps}: the ket is contracted from its two ends into its 16 amplitudes (steps 1 to 3), so is the bra
(4 to 6), and the last step joins them over the physical indices; for four sites this is the cheapest order:
\example{diagrampath}
\keepfor[6]{diagramsteps}\texttt{"Steps" -> True} gives one picture per step, with the pair contracted next highlighted, all with the same plot
range, so that \texttt{ListAnimate} shows the contraction. A ring of three tensors:
\example{diagramsteps}

\keepfor[8]{tree}\subsection{The contraction tree}
\fn{EinsumTree} draws the tree of the path: the leaves are the tensors, and each box is an intermediate, with its indices,
the indices summed at that step and its cost; the title gives the total cost, the naive one, the speed-up and the largest
intermediate. The matrix chain of section~\ref{sec:order}:
\example{tree}

% =====================================================================================================================
\keepfor[12]{efdefine}\section{Compiled einsums and the plan cache}
\fn{Einsum} keeps the plan of every call (the parsed specification, the path and the kernel of each step) in a cache,
keyed by the specification and the shapes: a second call with the same shapes goes straight to the kernels.
\fn{EinsumFunction} makes the plan once, for given shapes, and gives a function to be applied in a loop. The transfer
operator of an MPS, $x\mapsto\sum_sA^{s\,\mathsf T}xA^s$, that is $\sum_{ab}x_{ab}A_{asc}A_{bsd}$, has the summary of its
plan:
\example{efdefine}
Its properties are those of the path:
\example{efproperties}
Three hundred applications of the transfer operator, normalized, converge to its dominant eigenvector: the eigenvalue is
that of the $64\times64$ transfer matrix, built by another einsum:
\example{efloop}
\fn{EinsumClearCache[]} empties the cache.

% =====================================================================================================================
\section{Symbolic, exact, complex and sparse arrays}
The kernels are \texttt{Dot}, \texttt{Transpose}, \texttt{Total} and products of arrays, so the entries can be anything:
symbols, exact numbers, complex numbers. A product of symbolic matrices:
\example{symbolic}
The cross product, with the Levi-Civita tensor \texttt{LeviCivitaTensor[3]} (a \texttt{SparseArray}), and the identity
$\sum_i\epsilon_{ijk}\epsilon_{imn}=\delta_{jm}\delta_{kn}-\delta_{jn}\delta_{km}$:
\example{cross}
\example{epsilondelta}
The Pauli matrices as one $3\times2\times2$ array: $\sum_a\sigma^a\sigma^a=3$, and
$\sigma^a\sigma^b=\delta_{ab}+i\sum_c\epsilon_{abc}\sigma^c$:
\example{pauli}
\example{pauliproduct}
A \texttt{SparseArray} stays sparse in the steps that are matrix products; the other kernels make it dense. The square of
the hopping matrix of a ring of 1000 sites:
\example{sparse}

% =====================================================================================================================
\section{Physics}

\keepfor[10]{ising}\subsection{The partition function of the Ising model}
On a lattice, $Z=\sum_{\{s\}}\prod_{\langle ij\rangle}W_{s_is_j}$ with $W_{ss'}=e^{\beta Jss'}$: each bond is a $2\times2$
matrix, and each spin an index shared by all the bonds that touch it. On an open $3\times3$ lattice with the sites
$a,\dots,i$ the inner spins are in three or four bonds, the hyperedges of the diagram, and the spins at the corners, in
two bonds, are plain bonds:
\example{ising}
With $x=e^{\beta J}$, the partition function is a polynomial in $x$ and $1/x$, exact:
\example{isingz}
It is the sum over the $2^9$ configurations:
\example{isingcheck}

\subsection{Matrix product states, exactly}
The AKLT state of the spin-1 chain~\cite{AKLT87} is an MPS with $\chi=2$: $A^{+1}=\sqrt{2/3}\,\sigma^+$,
$A^0=-\sqrt{1/3}\,\sigma^z$, $A^{-1}=-\sqrt{2/3}\,\sigma^-$. On a ring of $L$ sites, \texttt{akltRing} writes its norm, a
trace over the bonds, and with operators on some sites an expectation value. With exact entries the results are exact:
\example{aklt}
They are $\mathrm{Tr}\,E^L=1+3(-1/3)^L$, from the eigenvalues of the transfer matrix $E$, $1$ and $-1/3$ three times:
\example{akltnorm}
The correlation $\langle S^z_1S^z_{1+r}\rangle$ on a ring of ten sites, for $r=1,\dots,4$; as the ring grows it tends to
its value on the infinite chain, $\frac43(-\frac13)^r$, that is $-\frac49$, $\frac4{27}$, $-\frac4{81}$, $\frac4{243}$:
\example{akltcorrelation}

\subsection{A Kubo formula}
The optical conductivity of a two-band chain, the Rice--Mele model~\cite{RM82}
$H(k)=(1+0.6\cos k)\,\sigma_x+0.6\sin k\,\sigma_y+0.4\,\sigma_z$, at half filling. With the band energies $E_n(k)$, the
occupations $f_n(k)$ and the current $J(k)=\partial_kH(k)$ in the band basis, $J_{nm}(k)=\langle nk|J(k)|mk\rangle$, the
Kubo formula reads
\[
\sigma(\omega)\propto\frac iN\sum_k\sum_{nm}J_{nm}(k)\,J_{mn}(k)\,\frac{f_n(k)-f_m(k)}{E_m(k)-E_n(k)}\,
\frac1{\omega-[E_m(k)-E_n(k)]+i\eta}.
\]
The bands on $N=400$ values of $k$:
\example{kubomodel}
The currents in the band basis, $J_{nm}(k)=\sum_{ij}\bar u_{ni}(k)J_{ij}(k)u_{mj}(k)$, for all $k$ in one einsum, where $k$
is in three tensors:
\example{kubocurrents}
The last two factors of the formula are an array $W_{knm\omega}$, on 401 frequencies (zero for $n=m$):
\example{kuboweights}
\keepfor[8]{kubo}The formula is then one einsum, written as it reads: $k$, $n$ and $m$ each appear in three tensors. The absorption,
$\mathrm{Re}\,\sigma(\omega)$, starts at the gap, $2\min_k E_2(k)\approx1.13$, and ends at $2\max_kE_2(k)\approx3.30$,
with the peaks of the van Hove singularities at the edges of the band of transitions:
\example{kubo}
It is the sum written out:
\example{kubocheck}
The package computes it in two steps, the product of the currents and its contraction with the weights:
\example{kubopath}

\subsection{A loop of propagators}
With the Green's functions $G_k(z)=[z-H(k)]^{-1}$, the trace of a loop of three, $\mathrm{Tr}\,G_k(z_1)G_k(z_2)G_k(z_3)$
for every $k$, as in a triangle diagram of nonlinear response, is \texttt{"kab,kbc,kca->k"}. In the band basis it is
$\sum_n\prod_i[z_i-E_n(k)]^{-1}$:
\example{propagators}

% =====================================================================================================================
\section{How a contraction runs}
Each pairwise contraction is classified by its indices: batch indices (in both operands and kept), contracted indices (in
both and summed) and free indices (in one operand). Before that, the indices that appear in one operand only and are not
needed afterwards are summed with \texttt{Total}, and diagonals are taken by one \texttt{Transpose} with repeated levels.
Then one kernel runs:
\begin{itemize}
\item \textbf{Dot}: no batch index. Each operand becomes a matrix with one \texttt{Transpose} and one \texttt{ArrayReshape}
  at most, none when its indices are already in order, and the product is one call of BLAS. Among the orders of the
  operands and of their indices, the one that moves the least data, the final result included, is chosen.
\item \textbf{Batch}: batch indices, $C_{bik}=\sum_jL_{bij}R_{bjk}$. A cost model chooses between a loop of \texttt{Dot}
  (large matrices) and a vectorised product with \texttt{Total} on packed arrays (many small matrices).
\item \textbf{Reduce}: two operands with the same indices, as \texttt{"kh,kh->h"}: the product in the layout of one of them
  and \texttt{Total} over the contracted levels, instead of transposing both.
\item \textbf{Hadamard}, \textbf{Prefix}, \textbf{Suffix}: no summed index; the elementwise product, or a broadcasting
  product when the indices of one operand are the first or the last of the other. \textbf{Scale}: a scalar operand.
\end{itemize}
The plan of an \fn{EinsumFunction} tells the kernel of each step (and \texttt{Dot} or \texttt{Vec}, the vectorised product,
for a batch):
\example{kernels}
The flops count the arithmetic; the kernels decide how fast it goes. A \texttt{Dot} runs at the speed of BLAS, and a batch of
$2\cdot10^4$ products of $4\times4$ matrices runs as a few operations on packed arrays, not as $2\cdot10^4$ calls.
Intermediates are freed as soon as they are used. \texttt{Examples/Benchmarks.wl} measures the timings on your computer,
against \texttt{Dot}, \texttt{MapThread}, \texttt{TensorContract} and \texttt{Total}.

% =====================================================================================================================
\section{How it is checked}
The test suite (\texttt{wolfram -script Tests/RunTests.wls}, 735 tests) checks 213 einsums against \texttt{numpy.einsum},
computed in exact integer arithmetic: 63 chosen to cover every kernel, repeated with every path strategy and both batch
methods, and 150 random ones. On 40 networks it checks that the cost of the \opt{Optimal} path equals that of an exhaustive
search over all the orders, and that \opt{Greedy} is never cheaper. It compares results in machine precision, real and
complex, with sums written out, and checks symbolic and sparse inputs, the syntaxes, the messages, \fn{EinsumFunction},
the diagrams and the update check. The package was first written without Mathematica and checked with Mathics3, an
open-source interpreter of the Wolfram Language, against \texttt{numpy.einsum} on about 2500 einsums, with every path and
every kernel.

The exhaustive search is easy to repeat: the five tensors of section~\ref{sec:order} can be contracted in
$10\cdot6\cdot3=180$ orders, and the cheapest one costs as much as the one the package finds:
\example{exhaustive}
And a ring of three tensors, in integers, against its sum written out:
\example{explicitsum}

\begin{thebibliography}{9}\small
\bibitem{SG18} D.~G.~A. Smith and J.~Gray, J.\ Open Source Softw.\ {\bf 3}, 753 (2018).
\bibitem{GKN04} E.~R. Gansner, Y.~Koren, and S.~North, in \emph{Graph Drawing 2004}, Lect.\ Notes Comput.\ Sci.\ {\bf 3383},
239 (Springer, 2005).
\bibitem{AKLT87} I.~Affleck, T.~Kennedy, E.~H. Lieb, and H.~Tasaki, Phys.\ Rev.\ Lett.\ {\bf 59}, 799 (1987).
\bibitem{RM82} M.~J. Rice and E.~J. Mele, Phys.\ Rev.\ Lett.\ {\bf 49}, 1455 (1982).
\end{thebibliography}
\end{document}
